\scanpage{42}
% archive sheet number: 41
\[
E(i,j)=E/(\sigma_0,\ldots,\sigma_{i-1};\sigma_{j+1},\ldots,\sigma_n)
=E/G^-(i)\times G^+(j)=D_{i,i+1,\ldots,j}
\]
et l'application
\[
E(i,j)\longrightarrow D_{i,j}\ \deleted{\unsure{D_j}}
\]
\margin{$0\leq i\leq i+2\leq j\leq\unsure{n+1}$}

Le groupe engendré par $\sigma_{i+1},\ldots,\sigma_{j-1}$ opère sur $E(i,j)$, avec comme quotient $D_{i,j}$. Sans doute il opère \uline{cartographiquement}, et on trouve une $(j-i-2)$-carte, dont les composantes connexes correspondent aux $d_{ij}\in D_{ij}$; \inserted{sur $\Pi(d_{ij})$} les facettes de dim $h$ $(0\leq h\leq j-i-2)$ correspondent aux drapeaux de type $(i,\underbrace{i+h+1}_{\inserted{i+1\leq\ \leq j-1}},j)$ au dessus du drapeau $d_{ij}$.

Cas particuliers:
\[
\begin{aligned}
j=i+2&\quad\text{\unsure{un peu idiot}}\\
j=i+3&:\quad\text{on trouve un polygône, d'où}\\
&\qquad\text{un \uline{invariant}}\quad\unsure{p|\rho}(d_{ij})\in\overline{\mathbb{N}}^*
\end{aligned}
\]
\[
(\text{où }1\leq p(d_{ij})\leq+\infty)
\]

\deleted{On \unsure{s'intéresse au}}
\deleted{Introduisons}
Formellement les $\Pi^{\unsure{-|=}}(f_i)$ correspondent à \deleted{\unsure{au}} $\Pi(d_{-1,i\deleted{\unsure{f_i}}})$, $d_{-1,i}$ défini par $f_i$ \deleted{\unsure{pour $\Pi$}} \deleted{\unsure{$(-1)$-facettes}} et $\Pi^+(\unsure{f^i|f_i|f'})$ \deleted{par} à $\Pi(d_{i,n+1})$, où $d_{i,n+1}$ est défini par $f_i$. Les invariants numériques
